\section{Future Work}\label{sec:future}

This appendix outlines the development roadmap beyond the current submission, organised by feasibility so that readers can distinguish between improvements achievable with existing hardware and those requiring additional investment or regulatory approval.

The system presented in this report demonstrates a complete autonomous SAR pipeline from search through detection to landing, but several directions offer clear paths toward operational capability. The following subsections organise future work by feasibility, distinguishing near-term enhancements achievable with the existing hardware from longer-term developments that would require additional equipment, regulatory approval, or multi-platform coordination.

\subsection{Inference Acceleration}\label{sec:future-inference}

The TFLite backend achieves 4.8\,FPS on the Raspberry Pi~5, which is adequate for the current search speed but leaves limited margin for faster flight or more complex models. Four acceleration strategies are identified in order of increasing implementation effort:

\begin{enumerate}
    \item \textbf{FP16 XNNPACK.} The Pi~5's Cortex-A76 cores support native half-precision arithmetic. Enabling FP16 execution in the XNNPACK delegate requires no model changes and is expected to approximately double throughput to ${\sim}10$~FPS, based on ARM's published benchmarks for similar workloads~\cite{david2021tensorflow}.
    \item \textbf{Threaded pipeline.} The current architecture serialises frame capture and inference. Overlapping these stages in a producer--consumer arrangement would reduce effective per-frame latency by 20--30\,\%, as the camera sensor and CPU operate on independent hardware. The Python \texttt{threading} module is sufficient because the Global Interpreter Lock (GIL) is released during both camera I/O and TFLite native inference calls.
    \item \textbf{NCNN backend.} An NCNN export of the YOLOv8n model has already been generated and is stored in the repository (\texttt{cv\_models/sar\_v2\_1088/ncnn/}). NCNN is optimised for ARM NEON instructions and is expected to reach ${\sim}15$~FPS on the Pi~5~\cite{ncnn}. Integration requires replacing the TFLite inference call in \texttt{vision.py} with an NCNN loader; the detection interface (\texttt{detect\_in\_image} $\rightarrow$ found, $x$, $y$, confidence) remains unchanged.
    \item \textbf{Hailo-8L NPU.} The Raspberry Pi AI~HAT+ (\pounds55) provides a dedicated neural processing unit capable of 80+ FPS for YOLOv8n~\cite{hailo2023datasheet}, removing inference as a bottleneck entirely. This would enable real-time detection at full camera frame rate and free CPU cycles for concurrent tasks such as visual odometry or on-board path replanning.
\end{enumerate}

Of these, the first two are immediately achievable with no additional hardware and represent the recommended next steps. The NCNN backend requires validation on the Pi to confirm that the exported model reproduces detections equivalent to the TFLite baseline. The Hailo accelerator, while transformative, introduces a hardware dependency and a model conversion workflow (HEF format compiled via an x86 Linux toolchain) that places it in the medium-term category.

\subsection{Detection Improvements}\label{sec:future-detection}

The YOLOv8n model achieves $\text{mAP}_{50} = 0.995$ on the validation set, but operational robustness under varied field conditions requires further work:

\begin{itemize}
    \item \textbf{Temporal voting.} Requiring a detection to persist in at least $N$ of $M$ consecutive frames before triggering investigation would suppress transient false positives caused by shadows, debris, or image artefacts. The spatial clustering system already provides partial temporal consistency, but an explicit frame-count threshold would add a complementary filtering layer at negligible computational cost.
    \item \textbf{Size and aspect-ratio filtering.} At a given altitude, a human casualty occupies a predictable range of bounding-box dimensions and aspect ratios. Rejecting detections outside this range---for example, boxes that are too small (sensor noise) or too wide (vehicles, paths)---would reduce the false-positive rate without retraining the model.
    \item \textbf{Hard negative mining.} The current training set includes 50 negative images (frames with no target). Systematically collecting false positive crops from flight footage and adding them as hard negatives during retraining would teach the model to distinguish the dummy from visually similar background features such as rocks, litter, or trail markers.
    \item \textbf{Expanded real training data.} Each flight produces thousands of frames that can be labelled with the existing annotation tool. Incorporating 200+ real images spanning varied lighting, altitude, pose, and background conditions would reduce the synthetic-to-real domain gap that currently limits generalisation~\cite{tremblay2018training}.
\end{itemize}

\subsection{Multi-Target Support}\label{sec:future-multi}

The simulator already implements spatial clustering and priority scoring for multiple simultaneous detections. Extending the flight system to support a multi-target queue would allow the drone to detect, classify, and record the positions of several casualties in a single sortie, revisiting unconfirmed detections after completing the primary search pattern. The ground station interface already supports per-target classification (confirmed, interest, false positive), so the principal remaining work lies in the state machine's transition logic for queueing and revisiting targets. In mass-casualty scenarios, this capability would enable a single sortie to locate and prioritise multiple victims before ground teams are dispatched.

\subsection{Thermal Imaging}\label{sec:future-thermal}

Visual detection is fundamentally limited by lighting conditions and the contrast between the casualty and the surrounding terrain. A thermal camera operating in the long-wave infrared band (8--14\,$\mu$m) would detect body-heat signatures independently of visible appearance, enabling night-time and low-visibility operations~\cite{rudol2008human}. Thermal and visible imagery could be fused at the decision level: a detection confirmed in both modalities would carry substantially higher confidence than either alone. Lightweight thermal modules compatible with the Raspberry Pi (e.g., FLIR Lepton~3.5) are available for under \pounds200 and produce $160\times120$ images sufficient for person-scale detection from 30\,m altitude. However, thermal imaging introduces additional challenges---sun-heated objects produce signatures similar to body heat, and wind chill reduces the thermal contrast between a casualty and the ambient background~\cite{thermal_sar}---so it complements rather than replaces the visual pipeline.

\subsection{Communication}\label{sec:future-comms}

The current system relies on a direct MAVLink telemetry link between the ground station laptop and the aircraft, limiting operational range to the Wi-Fi or radio link budget (typically 500\,m to 1\,km). Two enhancements would extend this range:

\begin{itemize}
    \item \textbf{4G/LTE telemetry.} A cellular modem on the aircraft would provide telemetry and command links over the mobile network, extending range to wherever cellular coverage exists. This is particularly relevant for rural SAR operations where the search area may extend several kilometres from the operator. Lightweight LTE modules (e.g., Sixfab 4G HAT) weigh under 30\,g and integrate directly with the Pi~5. Typical latency (50--150\,ms) is acceptable for supervisory control but would require a local autonomy layer to handle time-critical decisions (e.g., obstacle avoidance) without round-trip delay.
    \item \textbf{Real-time video to command centre.} Streaming compressed video (H.264 over WebRTC or RTSP) to a remote incident command post would allow SAR coordinators to observe detections as they occur, rather than relying on post-flight review. The Pi~5's hardware video encoder can produce a 720p stream at minimal CPU cost, but upstream bandwidth constraints on cellular networks (1--5\,Mbps typical) would require adaptive bitrate control.
\end{itemize}

Both enhancements depend on obtaining Beyond Visual Line of Sight (BVLOS) operational authorisation from the UK Civil Aviation Authority under CAP~722~\cite{cap722}, which requires a detect-and-avoid capability for other airspace users, redundant communication links, and a comprehensive safety case.

\subsection{Swarm Operations}\label{sec:future-swarm}

Deploying multiple drones to divide a search area would reduce coverage time approximately linearly with fleet size. Each aircraft would fly an independent sub-region while sharing detection events over a mesh network, enabling collaborative coverage planning that avoids redundant passes~\cite{waharte2010supporting}. The lawnmower pattern generator already accepts arbitrary polygons and could partition a large search area into sub-polygons assigned to individual aircraft. This is the most aspirational direction discussed here: the primary challenges are inter-drone communication reliability, deconfliction of overlapping flight paths, and a centralised operator interface that aggregates detections from all platforms without overwhelming the operator~\cite{multiuav}. Nevertheless, the modular software architecture---where detection, planning, and control are independent modules communicating through well-defined interfaces---provides a foundation that would support multi-agent extension without fundamental redesign.
